\section{SYNTHETIC MECHANISM DETAILS}
\noindent The maintenance environment exposes binary ventilation faults through two independent noisy clues. Faults have equal prior probability, with clue accuracies 0.8 and 0.6. An agent chooses a filter replacement or sensor reset from a public manual. Two fresh model samples are drawn per job. The first action executes; the second supplies agreement information. Review may correct the executed action and provide definitive feedback for later jobs.\draftcredit

An initially wrong action incurs downtime rate $c_i$ until correction or deadline, plus terminal incorrect-closure penalty $K_i$. A timely correction at $T$ prevents
\begin{equation}
 G_i(T)=\ind[T\leq d_i]\{c_i(d_i-T)+K_i\}.
 \label{eq:maintenance}
\end{equation}
The original workload uses three agents, two jobs each and a 12-tick horizon. Releases occur at wave bases 0 and 6 with zero-or-one jitter. Relative windows are drawn from 1, 2 and 5, downtime rates from 0, 1 and 3, and terminal penalties from 0 and 8.

A constructed competition workload releases all three jobs together at 0 and 6. Each wave independently permutes windows 2, 4 and 5 and penalties 4, 8 and 12, with zero downtime. It was motivated by limited competition in earlier development, then evaluated on new frozen scenarios. A larger workload has three or six agents, three jobs each, a 24-tick horizon and waves at 0, 8 and 16 with independent zero-or-one jitter. Windows are 2, 4 and 6, downtime rates 0, 1 and 3, and penalties 0, 4, 8 and 12. Stable per-agent streams pair the first three agents across team sizes. Public assignments are independent of hidden faults.

\subsection{Risk and Frozen Evaluation}
The primary risk estimate is an agreement-bin estimator fit on an earlier development set. Laplace smoothing gives $(e+1)/(n+2)$ for $e$ errors among $n$ jobs; bins with fewer than ten examples use pooled calibration risk. Initial first-action labels are included regardless of review selection. The frozen values are $13/89$ for agreement and pooled $18/98$ for sparse disagreement.

Alternative estimates use pooled risk or a public analytical reference with knowledge of the synthetic observation model. Following agreeing clues has error probability $1/7$; following the stronger disagreeing clue has error probability $3/11$. Choosing the opposite action gives the complementary probability. The reference does not inspect the realized fault. Both samples are retained in every condition. Prediction rescoring uses saved outputs; policy comparisons under alternative risk use actual new trajectories.

The frozen evaluation contains 2,752 episodes. Original and competition workloads each use 64 scenarios and six policies at one- and two-tick reviews. Alternative-risk comparisons use the first 32 scenarios of each workload for greedy and search. Larger-workload comparisons use 32 scenarios with four policies across both team sizes and capacities. A separately declared objective study uses 16 scenarios. The development and evaluation partitions remain separate.

\subsection{Results and Boundary Conditions}
Table~\ref{tab:maintenance} reports central search--greedy contrasts. At two ticks, search gains in constructed competition but has little measured advantage in the original workload. It wins, ties and loses against greedy on 15, 42 and 7 competition scenarios, versus 3, 59 and 2 original scenarios. Search versus EDF in two-tick competition is $-2.5625$ points, with interval [$-4.0000,-1.3125$]. At one tick, search and EDF both attain zero competition loss despite different review sequences on 62 of 64 scenarios.

\begin{table*}[t]
\caption{Frozen maintenance results. Differences are search minus greedy mean loss with 95\% paired scenario-bootstrap intervals. Workload distributions are analyzed separately.}
\label{tab:maintenance}\centering\small
\begin{tabular}{llrrrrl}
\toprule
Workload & Agents / scenarios & Review ticks & Search & Greedy & EDF & Paired difference \\
\midrule
Original & 3 / 64 & 2 & 7.328 & 7.422 & 7.656 & $-0.094\;[-0.563,\;0.344]$ \\
Competition & 3 / 64 & 2 & 1.688 & 3.000 & 4.250 & $-1.313\;[-2.250,-0.438]$ \\
Competition & 3 / 64 & 1 & 0.000 & 0.375 & 0.000 & $-0.375\;[-0.688,-0.125]$ \\
Larger & 6 / 32 & 1 & 16.188 & 22.375 & 23.688 & $-6.188\;[-9.376,-3.281]$ \\
Larger & 6 / 32 & 2 & 36.563 & 36.781 & 46.000 & $-0.219\;[-2.313,\;1.969]$ \\
\bottomrule
\end{tabular}
\end{table*}

The larger workload gives a stronger one-tick planning benefit and a small, uncertain two-tick difference. Reduced capacity does not monotonically increase the value of planning. Some requests are already infeasible when the reviewer becomes free. In the original two-tick search runs, only 32 of 156 dispatches have multiple eligible requests, whereas competition has 224 of 256. Such public contention is a prerequisite for many sequencing effects, but it does not guarantee competing actual errors.

Risk calibration and allocation utility diverge. The frozen estimator predicts roughly 0.15 error while observed errors reach about 0.25--0.32. On two-tick search trajectories, analytical risk improves Brier score from 0.213 to 0.157 in the original workload and from 0.200 to 0.164 in competition. Yet fresh matched-prefix alternative-risk rollouts yield no consistent allocation gain. Better probability scores alone do not establish better counterfactual policy performance.

Weighted cost also differs from task correctness. In the six-agent one-tick workload, search has lower loss than EDF but more incorrect closures. Adding a declared penalty $\lambda$ per incorrect closure to the search objective changes both intervention benefit and eligibility. With $\lambda=8$ rather than zero, original cost increases by 1.75 points [$0.188,3.438$], while incorrect closures decrease by 0.75 [$0.375,1.125$] per scenario. The common augmented objective $\loss+8U$ improves by 4.25 points, where $U$ counts incorrect closures. At two ticks this improvement weakens and the augmented objective is slightly worse. This is a tradeoff between objectives, not a universally superior policy.

